\documentclass[UTF8,12pt,a4paper]{ctexart}

\usepackage{amsmath}
\usepackage{amssymb}
\usepackage{graphicx}
\usepackage{enumerate}
\usepackage[margin=1in]{geometry}
\geometry{a4paper}
\usepackage{booktabs}
\usepackage{array}
\usepackage{longtable}
\newcolumntype{L}[1]{>{\raggedright\arraybackslash}p{#1}}
\newcolumntype{M}[1]{>{\centering\arraybackslash}p{#1}}
\usepackage{fancyhdr}
\pagestyle{fancy}
\fancyhf{}
\usepackage{xcolor}




% ------------------- 设置超链接，便于跳转 -----------------
\usepackage{enumitem}
\usepackage{hyperref}
\hypersetup{
    pdfborder={0 0 0},    % 消除超链接边框
    colorlinks=true,       % 将边框改为颜色标记（可选）
    linkcolor=black,       % 内部链接颜色（目录、引用等）
    citecolor=black,       % 引用颜色
    urlcolor=blue          % URL链接颜色（保持蓝色以便区分）
}




% ------------------------ 标题区 ----------------------------
\title{HW3}
\author{
	姓名：\underline{冯峻}~~~~~~
	学号：\underline{523031910148}~~~~~~}
\date{\today}
\pagenumbering{arabic}

\begin{document}



% ------------------------ 页眉和页脚 ----------------------------
\fancyhead[L]{冯峻}
\fancyhead[C]{Reisman模型}
\fancyfoot[C]{\thepage}

\maketitle

\section{问题背景}

\subsection{初始快速氧化现象}

在硅的热氧化过程中，经典的Deal-Grove模型能够很好地描述较厚氧化层（通常 > 30 nm）的生长动力学。然而，在氧化的初始阶段，特别是当氧化层厚度小于30 nm时，实验观察到的氧化速率明显快于Deal-Grove模型的预测。这种现象被称为\textbf{初始快速氧化（Initial Rapid Oxidation）}。

\section{Reisman et al. 模型}

\subsection{模型核心思想}

\subsubsection{基本假设}

Reisman等人提出的模型基于以下关键假设：

\begin{enumerate}
    \item \textbf{双层氧化层结构}
  
    在氧化初期，氧化层由两个不同性质的层组成：(1) 靠近Si基底的\textbf{界面层}（Interface Layer）；(2) 靠近气相的\textbf{体层}（Bulk Layer）
    
    界面层相对体层有不同的性质，例如密度、折射率和应力状态与体层不同、对氧化剂的扩散系数更大、厚度相对固定，约为几纳米。

    \item \textbf{氧化速率的控制机制}
   
    在初始阶段，氧化主要通过界面层进行，由于界面层的高扩散系数，氧化速率较快，随着体层厚度增加，扩散路径增长，速率逐渐降低。
    
\end{enumerate}

\subsubsection{数学描述}

\textbf{氧化层厚度的组成：}

总氧化层厚度 x 可以表示为：
\begin{equation}
x = x_i + x_b
\end{equation}

其中，$x_i$代表界面层厚度（相对恒定），$x_b$代表体层厚度（随时间增长）。

\textbf{修正的氧化动力学方程：}

Reisman模型对Deal-Grove方程进行修正，考虑界面层的影响：

对于界面层：
\begin{equation}
F_i = \frac{D_i C^*}{x_i}
\end{equation}

对于体层：
\begin{equation}
F_b = \frac{D_b C^*}{x_b}
\end{equation}

其中：
\begin{itemize}
    \item $D_i$：界面层中氧化剂的有效扩散系数
    \item $D_b$：体层中氧化剂的有效扩散系数
    \item $D_i \gg D_b$（界面层扩散系数远大于体层）
\end{itemize}

氧化剂通过双层结构的总通量为：
\begin{equation}
\frac{1}{F_{total}} = \frac{1}{F_i} + \frac{1}{F_b} = \frac{x_i}{D_i C^*} + \frac{x_b}{D_b C^*}
\end{equation}

因此：
\begin{equation}
F_{total} = \frac{C^*}{\frac{x_i}{D_i} + \frac{x_b}{D_b}}
\end{equation}

\subsubsection{初始快速氧化的解释}

\textbf{1. 在氧化初期（$x_b \ll x_i$或$x_b$很小）：}

当体层厚度很小时，主要的扩散阻力来自于界面层：
\begin{equation}
F_{total} \approx \frac{D_i C^*}{x_i}
\end{equation}

由于$D_i \gg D_b$，此时的氧化通量远大于Deal-Grove模型预测的值，导致\textbf{快速氧化}。

\textbf{2. 在氧化后期（$x_b \gg x_i$）：}

当体层厚度远大于界面层厚度时：
\begin{equation}
F_{total} \approx \frac{D_b C^*}{x_b}
\end{equation}

此时扩散主要由体层控制，氧化速率回归到Deal-Grove模型的预测，即：
\begin{equation}
x^2 \approx B t
\end{equation}

这样就实现了从初始快速氧化到Deal-Grove模型描述的平滑过渡。

\section{模型的物理机制分析}

\subsection{界面层的微观结构}

\subsubsection{结构特征}

界面层的特殊性质可能源于以下微观结构：

\textbf{过渡层性质}
    
Si/SiO$_2$界面不是突变的，而是存在一个过渡区，过渡区中Si-O键的配位数和键长逐渐变化，由晶态Si的四配位逐渐过渡到非晶SiO$_2$的随机网络结构

\textbf{应力梯度}
    
氧化过程伴随体积膨胀（Si → SiO$_2$），界面附近存在较大的压应力，应力可能导致氧化层结构的局部畸变，畸变的结构可能为氧化剂扩散提供更多通道。
    


\subsubsection{高扩散系数的原因}

界面层具有高扩散系数（$D_i \gg D_b$）的可能原因：

\begin{itemize}
    \item \textbf{结构无序度：}界面区域的无序度更高，提供了更多的扩散路径
    \item \textbf{应力辅助扩散：}压应力可能降低了扩散激活能
    \item \textbf{缺陷辅助扩散：}点缺陷和悬挂键作为扩散的快速通道
    \item \textbf{键角变化：}Si-O-Si键角的分布范围更宽，利于分子穿行
\end{itemize}

\end{document}
